Quando um computador inicia uma aplicação, o sistema operacional aloca 
automaticamente a memória necessária para a sua execução em diferentes 
segmentos, cada um com propriedades e funções específicas.

\begin{figure}[htbp]
\centering
\includegraphics[scale=0.2]{memory-segments}
\label{fig:memorySegmentsDiagram}
\end{figure}

\subsection{Text}

A seção \textbf{Text}, também conhecida como seção de \textbf{Texto}, contém o 
código executável do programa. As instruções de máquina presentes nesta seção 
são acessadas pelo processador para execução da aplicação. Importante destacar 
que esta seção é somente de leitura, o que significa que, uma vez definida, seu 
conteúdo não pode ser modificado durante a execução do programa. Isso é 
essencial para a integridade e segurança do código executável.

\subsection{Data}

A seção \textbf{Data}, também conhecida como seção de \textbf{Dados 
Inicializados}, é onde são armazenadas as variáveis globais e estáticas do 
programa que foram inicializadas com valores específicos. É importante 
destacar que a seção de dados inicializados não é apenas de leitura; ela pode 
ser modificada durante a execução do programa.

Considere o seguinte código:

\begin{code}
    static char *z = "Eu sou Z\n"
    char *w = "Eu sou W\n"

    void func(void)
    {
        static const char *y = "Eu sou Y\n"
    }

    int main(void)
    {
        const char *x = "Eu sou X\n"
        func();

        return 0;
    }
\end{code}

No código acima, são definidas as variáveis \texttt{x}, \texttt{y}, \texttt{z} 
e \texttt{w}. Todas elas, exceto \texttt{x}, são direcionadas para a seção 
\textbf{Data}. Isso ocorre porque as variáveis globais, independentemente de 
usarem a diretiva \textbf{static}, e as variáveis locais que utilizam a 
diretiva \textbf{static}, são alocadas na seção \textbf{Data} caso sejam 
inicializadas corretamente.

\subsection{BSS}

A seção \textbf{BSS}, também conhecida como seção de \textbf{Dados Não 
Inicializados}, armazena variáveis globais e estáticas que não foram 
inicializadas explicitamente pelo programador. Frequentemente, a seção 
\textbf{BSS} é considerada juntamente com a seção \textbf{Data} por ambas 
conterem variáveis globais e estáticas, diferenciando-se apenas pelo estado 
de inicialização das variáveis.

Repetindo o mesmo código da seção anterior de forma ilustrativa:

\begin{code}
    static char *z;
    char *w;

    void func(void)
    {
        static const char *y;
    }

    int main(void)
    {
        const char *x;
        func();

        return 0;
    }
\end{code}

As variáveis \texttt{y}, \texttt{z} e \texttt{w}, sendo declaradas como globais 
ou estáticas sem uma inicialização explícita, são alocadas na seção \textbf{BSS} 
e inicializadas automaticamente com zero. Isso é uma característica padrão do 
compilador para assegurar que todas as variáveis na seção \textbf{BSS} comecem 
com um estado previsível.

Por outro lado, a variável \texttt{x} apesar de não inicializada é declarada 
localmente, portanto não pertence à seção \textbf{BSS}. Em C, variáveis locais 
não inicializadas como \texttt{x} podem conter valores indeterminados, pois o 
compilador não realiza nenhuma inicialização automática para elas. Isso implica 
que \texttt{x} pode conter qualquer valor aleatório presente na stack no 
momento de sua declaração.

\begin{ps}
    \textbf{BSS} é uma abreviação para: \textbf{Block Started By Symbol}. 
\end{ps}

\subsection{Heap}

A seção \textbf{Heap} é utilizada para armazenar dados alocados e liberados dinamicamente pelo programa durante sua execução, por meio de funções como \texttt{malloc}, \texttt{calloc}, \texttt{realloc} e \texttt{free}.

É importante destacar que, ao contrário das outras seções de memória, a seção \textbf{Heap} não é necessariamente contígua na memória. Ela pode utilizar técnicas de mapeamento e alocação não contígua para ajustar o tamanho dos blocos de memória alocados dinamicamente. Devido a essa maior complexidade e dinamicidade dos blocos de memória na \textbf{Heap}, acessar variáveis localizadas nessa seção pode ser mais lento em comparação com as seções estáticas.

Além disso, a seção \textbf{Heap}, conforme ilustrado na figura \ref{fig:memorySegmentsDiagram}, começa imediatamente após a seção \textbf{BSS} e pode crescer em direção à seção \textbf{Stack} à medida que mais memória dinâmica é alocada e utilizada.

%TODO: Elaborar uma melhor explicação com mais exemplos do que é a Stack
\subsection{Stack}

A seção de memória da Stack é parte crucial da memória de um programa em C, 
responsável pelo armazenamento de variáveis locais e pela manutenção do escopo 
de funções.

A Stack, também conhecida como pilha, é uma estrutura de dados de natureza 
LIFO\footnote{Last In, First Out}, o que significa que o último item a ser 
inserido é o primeiro a ser retirado. No contexto da execução de um programa em 
C, a \textbf{Stack} é utilizada para armazenar variáveis locais, parâmetros de 
função e endereços de retorno.

Quando uma função é chamada, um novo frame de execução é criado na \textbf{Stack} 
para armazenar suas variáveis locais e outros dados relacionados. Esse frame é 
desalocado quando a função retorna, liberando assim a memória alocada para suas 
variáveis locais.

É importante ressaltar que a \textbf{Stack} possui um tamanho limitado, 
determinado pelo sistema operacional ou pelo compilador. Portanto, a alocação 
excessiva de memória na \textbf{Stack} pode levar a um estouro de 
pilha\footnote{também conhecido como Stack Overflow}, resultando em falha na 
execução do programa. Além disso, como a memória na \footnote{Stack} é alocada 
de forma estática durante a compilação, não é possível alocar ou liberar memória 
de forma dinâmica, como acontece na \textbf{Heap}.

